\documentclass[11pt,a4paper]{article}
\usepackage[utf8]{inputenc}
\usepackage[T1]{fontenc}
\ifdefined\pdfgentounicode
\input{glyphtounicode}
\pdfgentounicode=1
\fi
\usepackage{amsmath,amssymb,amsthm,mathtools}
\usepackage{geometry}
\geometry{margin=1in}
\usepackage{booktabs}
\usepackage{microtype}
\usepackage{cite}
\usepackage{url}
\def\UrlBreaks{\do\/\do-\do\_}
\usepackage{xcolor}
\usepackage{hyperref}

\hypersetup{
    colorlinks=true,
    linkcolor=blue,
    citecolor=red,
    urlcolor=blue,
    pdftitle={Certified Lower Bounds and Variational Dynamics for the Erdos Minimum Overlap Problem},
    pdfauthor={Jayden Xing},
    pdfsubject={Erdos Problem 36; dual certificates and interval arithmetic}
}

\newtheorem{theorem}{Theorem}[section]
\newtheorem{lemma}[theorem]{Lemma}
\newtheorem{proposition}[theorem]{Proposition}
\newtheorem{corollary}[theorem]{Corollary}
\theoremstyle{definition}
\newtheorem{definition}[theorem]{Definition}
\theoremstyle{remark}
\newtheorem{remark}[theorem]{Remark}

\newcommand{\R}{\mathbb{R}}
\newcommand{\Z}{\mathbb{Z}}
\newcommand{\N}{\mathbb{N}}
\newcommand{\Q}{\mathbb{Q}}
\newcommand{\E}{\mathbb{E}}
\newcommand{\supp}{\operatorname{supp}}
\newcommand{\sinc}{\operatorname{sinc}}

\title{Certified Lower Bounds and Variational Dynamics for the Erd\H{o}s Minimum Overlap Problem}

\author{
  Jayden Xing\thanks{Kreflux and FreezeSol. Email: \texttt{jayden@kreflux.com}. Code, certificates and Lean sources: \url{https://github.com/Destructive-Entity/erdos-minimum-overlap}.}
}

\date{November 2026}

\begin{document}

\maketitle

\begin{abstract}
Erd\H{o}s Problem \#36, the minimum overlap problem (1955), asks for the constant $\mu = \lim_{N \to \infty} M(N)/N$, where $M(N)$ is the minimum, over partitions $A \sqcup B = \{1, \dots, 2N\}$ with $|A| = |B| = N$, of the maximum over $k$ of the number of solutions of $a - b = k$ with $a \in A$ and $b \in B$. Working within the mean-conditioned dual certificate framework of Price, which builds on the Fourier-analytic approach of White, we prove $\mu > 0.38056171$ ($D_{\max} \le 2.62769470$). This improves the certified bounds of Price ($0.38055470$), Raval ($0.38055925$) and Deng ($0.38056070$) obtained in the same framework; higher bounds, $0.3805634$ and $0.38058706$, have been announced independently by Machillari Dowgu\~{n}iec and by Budala. The certificate uses 178 mean bins, with central bins $[-1/512,0]$ and $[0,1/512]$, White's weighted Parseval constraint on 400 harmonics, and off-grid phase-shifted cosine (line-bundle) constraints obtained from a complex lens enclosure of the Fourier transform. Every bin is verified in 160-bit Arb ball arithmetic, and the central window is cross-checked by an independent 192-bit MPFR verifier using a different certificate. We also prove an exact fourth-moment identity for the overlap function, record formal first-order conditions and a heuristic Duffing-type model for near-extremal profiles, and describe a Lean 4 file in which the elementary algebraic steps are checked and the analytic inputs and the numerical certificate are stated as explicit axioms.
\end{abstract}

\medskip
\noindent\textbf{Mathematics Subject Classification (2020).} 11B75; 42A05, 65G30.

\noindent\textbf{Keywords.} Minimum overlap problem, Erd\H{o}s problems, dual certificates, interval arithmetic.

\section{Introduction and Historical Context}

Let $N$ be a positive integer, and let $A \sqcup B = \{1, 2, \dots, 2N\}$ be a partition of the integers from $1$ to $2N$ into two disjoint subsets of equal cardinality $|A| = |B| = N$. For each integer lag $k \in \Z$, the intersection multiplicity between $A$ and the shifted set $B + k$ is
\begin{equation}
M_k(A, B) = |\{ (a, b) \in A \times B : a - b = k \} | = |A \cap (B + k)|.
\end{equation}
The minimum overlap problem of Erd\H{o}s \cite{erdos1955} (Problem C17 in Guy's collection \cite{guy2004} and Problem \#36 in \cite{erdosproblems36}) concerns the minimax quantity
\begin{equation}
M(N) = \min_{\substack{A \sqcup B = \{1, \dots, 2N\} \\ |A| = |B| = N}} \max_{k \in \Z} M_k(A, B)
\end{equation}
and the constant
\begin{equation}
\mu = \lim_{N \to \infty} \frac{M(N)}{N}.
\end{equation}

Erd\H{o}s observed that for even $N$ the choice $A = \{N/2+1, \dots, 3N/2\}$ gives $M(N) \le N/2$, and he initially conjectured that $\mu = 1/2$; Motzkin, Ralston and Selfridge \cite{motzkin1956} disproved this conjecture. A theorem of Swinnerton-Dyer, proved in Haugland's paper \cite{haugland1996}, shows that the limit exists and equals the value of a continuous minimax problem over measurable functions (Section~\ref{sec:continuous}). Previous work on $\mu$ includes the following.
\begin{enumerate}
\item \textbf{Early lower bounds.} An averaging argument gives $\mu \ge 1/4$. Scherk improved this to $\mu \ge 1 - 1/\sqrt{2} \approx 0.292893$ (see \cite{moser1959}), and Moser \cite{moser1959} proved $\mu \ge \sqrt{4 - \sqrt{15}} \approx 0.356394$. Moser and Murdeshwar \cite{mosermurdeshwar1966} introduced the continuous analogue studied below. Since a probability density bounded by $\mu$ has variance at least $1/(12\mu^2)$, the second-moment identity of Section~\ref{sec:fourth} gives the simpler bound $\mu \ge 1/\sqrt{8} \approx 0.353553$.
\item \textbf{Upper bounds.} Haugland constructed step functions showing $\mu < 0.382002$ \cite{haugland1996} and $\mu < 0.380926$ \cite{haugland2016}. Computer searches lowered the upper bound to $0.380924$ (AlphaEvolve \cite{alphaevolve2025}), $0.380876$ (TTT-Discover \cite{tttdiscover2026}) and $0.380868$ \cite{simpletes2026}; entries on the EinsteinArena leaderboard \cite{einsteinarena2026} have reached $0.38085858$ \cite{erdosproblems36}, and Budala \cite{budala2026} gives a construction with value below $0.3808566$.
\item \textbf{Convex programming.} White \cite{white2022, white2023} translated the problem into a family of convex programs in the Fourier coefficients and moments of the weight and proved $\mu \ge 0.379005$. Kim and Pilanci \cite{kimpilanci2026} added Toeplitz positivity constraints to White's program and obtained $\mu \ge 0.37912$.
\item \textbf{Mean-conditioned dual certificates.} Price \cite{price2026} partitioned the range $[-1, 1]$ of the mean $\bar{x}$ of the overlap function into 172 bins and certified a dual function for each bin in 160-bit Arb ball arithmetic, obtaining $\mu \ge 0.38055470$ ($D \le 2.62774313$). Raval \cite{raval2026} refined the central region and obtained $\mu > 0.38055925$, checked in Arb and MPFI. Deng \cite{deng2026} replaced the two central bins of Price by a certificate with 80 off-grid cosine constraints and White's weighted Parseval constraint on 400 harmonics, verified with a sixth-order Taylor MPFR verifier, and obtained $\mu > 0.38056070$.
\item \textbf{Later announcements.} Machillari Dowgu\~{n}iec \cite{machillari2026} added subadditivity constraints for the overlap profile and announced $\mu \ge 100000000/262768281 > 0.3805634$. Budala \cite{budala2026} announced $\mu \ge 0.38058706$ with exact-arithmetic checkers. Both bounds exceed the bound proved here.
\end{enumerate}

\subsection{Summary of Results in this Work}

\begin{enumerate}
\item \textbf{Certified lower bound.} Theorem~\ref{thm:main} states
\begin{equation}
\mu \ge \frac{1}{D_{\max}} > 0.38056171, \qquad D_{\max} \le 2.62769470.
\end{equation}
The certificate refines the central region to bins of width $1/512$ and combines mean, second-moment, weighted Parseval, cosine, sine and phase-shifted cosine constraints. It is verified on all 178 bins in 160-bit Arb ball arithmetic. An independent 192-bit MPFR verifier, applied to Deng's central certificate with the second-moment bound for the window $|\bar{x}| \le 1/1280$, gives $D \le 2.62769373$ on that window (Remark~\ref{rem:mpfr}).
\item \textbf{Fourth-moment identity.} Theorem~\ref{thm:fourth_moment} shows
\begin{equation}
\int_{-2}^2 x^4 M(x) \, dx = \frac{2}{5} + 4F_2 - 6F_2^2 + 8F_1F_3,
\end{equation}
so that the fourth moment $F_4$ of the weight cancels, and $\int_{-2}^2 x^4 M(x) \, dx \le \frac{16}{15} + |\E(M)|$.
\item \textbf{Variational structure.} Section~\ref{sec:variational} gives the centered decomposition of the overlap function, the vanishing of its boundary term for even profiles, formal first-order conditions for extremal profiles and a heuristic Duffing-type model. None of this is used in the proof of Theorem~\ref{thm:main}.
\item \textbf{Lean 4 file.} Section~\ref{sec:lean} describes \nolinkurl{formal/Erdos36.lean}, in which the elementary algebraic steps are checked and the analytic inputs, Swinnerton-Dyer's theorem and the numerical certificate are stated as axioms.
\end{enumerate}
Theorem~\ref{thm:main} improves the bounds of Price, Raval and Deng, but not the bounds announced by Machillari Dowgu\~{n}iec and Budala; it provides an independent certificate within the mean-conditioned framework.

\section{Continuous Formulation and Majorant Duality}
\label{sec:continuous}

\subsection{The Continuous Reduction of Swinnerton-Dyer}

Let $\mathcal{H}$ denote the class of measurable functions $f: [-1, 1] \to [0, 1]$ satisfying $\int_{-1}^1 f(t) \, dt = 1$. Extend $f$ by zero outside $[-1,1]$ and put $g = \mathbf{1}_{[-1,1]} - f$ on $\R$. The continuous overlap cross-correlation function $M_f: [-2, 2] \to [0, 1]$ is:
\begin{equation}
M_f(x) = \int_{\R} f(t) g(t+x) \, dt = \int_{-1}^1 f(t) \left(\mathbf{1}_{[-1, 1]}(t+x) - f(t+x)\right) \, dt.
\end{equation}
Because $f$ and $g$ are supported on $[-1, 1]$, the product $f(t)g(t+x)$ is non-zero only when $t \in [-1, 1]$ and $t+x \in [-1, 1]$, which forces $x = (t+x) - t \in [-2, 2]$. The function $M_f$ is nonnegative and supported on $[-2,2]$. Since $f,g \in L^2(\R)$, Cauchy--Schwarz gives $|M_f(x+h)-M_f(x)| \le \|f\|_2\|g(\cdot+h)-g\|_2$, which tends to zero as $h \to 0$ by continuity of translations in $L^2$. Thus $M_f$ is continuous. Fubini's theorem gives the mass identity:
\begin{equation}
\int_{-2}^2 M_f(x) \, dx = \left( \int_{-1}^1 f(t) \, dt \right) \left( \int_{-1}^1 g(s) \, ds \right) = 1.
\end{equation}

For $h \in L^1(\R)$, use the Fourier transform convention:
\begin{equation}
\widehat{h}(\xi) = \int_{\R} h(t) e^{-i\xi t} \, dt.
\end{equation}
Under this convention, for $f \in \mathcal{H}$ and $g(t) = 1 - f(t)$ supported on $[-1, 1]$, the transform of $\mathbf{1}_{[-1, 1]}$ is $\int_{-1}^1 e^{-i\xi t} \, dt = 2 \operatorname{sinc}(\xi)$, where $\operatorname{sinc}(\xi) = \sin(\xi)/\xi$ for $\xi \ne 0$ and $\operatorname{sinc}(0)=1$, so that $\widehat{g}(\xi) = 2\operatorname{sinc}(\xi) - \widehat{f}(\xi)$.

Writing $\widehat{f}(\xi) = u + iv$ and $s = \operatorname{sinc}(\xi)$, we have $\widehat{g}(\xi) = (2s - u) - iv$. The continuous overlap cross-correlation function $M_f(x) = \int_{\R} f(t)g(t+x) \, dt$ is supported on $[-2, 2]$ and satisfies $\widehat{M_f}(\xi) = \overline{\widehat{f}(\xi)} \, \widehat{g}(\xi)$. Expanding the product:
\begin{equation}
\widehat{M_f}(\xi) = (u - iv)((2s - u) - iv) = (2su - u^2 - v^2) - 2i s v.
\end{equation}
Therefore, $\operatorname{Re}\widehat{M_f}(\xi) = 2su - u^2 - v^2$ and $\operatorname{Im}\widehat{M_f}(\xi) = -2sv$. 
Since the Fourier transform decomposes as:
\begin{equation}
\widehat{M_f}(\xi) = \int_{-2}^2 M_f(x) \cos(\xi x) \, dx - i \int_{-2}^2 M_f(x) \sin(\xi x) \, dx,
\end{equation}
comparing real and imaginary parts yields:
\begin{align}
\int_{-2}^2 \cos(\xi x) M_f(x) \, dx &= \operatorname{Re}\widehat{M_f}(\xi) = 2su - u^2 - v^2, \\
\int_{-2}^2 \sin(\xi x) M_f(x) \, dx &= -\operatorname{Im}\widehat{M_f}(\xi) = 2sv.
\end{align}

By a theorem of Swinnerton-Dyer, proved in \cite{haugland1996} (see also \cite{white2022}), the limit $\mu$ exists and equals the continuous infimum over measurable functions:
\begin{equation}
\label{eq:sd}
\mu = \lim_{N \to \infty} \frac{M(N)}{N} = \inf_{f \in \mathcal{H}} \|M_f\|_\infty.
\end{equation}

\subsection{The Majorant Duality Principle}

To establish lower bounds on $\|M_f\|_\infty$, we use the following elementary duality principle, which underlies the certificates of \cite{price2026, deng2026}.
\begin{proposition}[Majorant Duality Principle]
\label{thm:majorant_duality}
Let $p: [-2,2] \to [0,\infty)$ be measurable and essentially bounded, with $\int_{-2}^2 p(x)\,dx=1$. Let $\mathcal{J}$ be finite, let each $\phi_j$ be a bounded measurable real-valued function, and suppose the real constants $B_j$ satisfy:
\begin{equation}
\int_{-2}^2 p(x) \phi_j(x) \, dx \le B_j, \quad \forall j \in \mathcal{J}.
\end{equation}
For any non-negative dual multipliers $\lambda_j \ge 0$, define the dual majorant:
\begin{equation}
q(x) = 1 + \sum_{j \in \mathcal{J}} \lambda_j (B_j - \phi_j(x)).
\end{equation}
Then:
\begin{equation}
\|p\|_\infty \ge \frac{1}{\int_{-2}^2 \max(0, q(x)) \, dx}.
\end{equation}
\end{proposition}

\begin{proof}
By non-negativity of $\lambda_j$ and the moment bounds $\int p \phi_j \le B_j$:
\begin{equation}
\int_{-2}^2 p(x) q(x) \, dx = \int_{-2}^2 p(x) \, dx + \sum_{j \in \mathcal{J}} \lambda_j \left( B_j - \int_{-2}^2 p(x) \phi_j(x) \, dx \right) \ge 1 + 0 = 1.
\end{equation}
For almost every $x$, nonnegativity gives:
\begin{equation}
p(x) q(x) \le p(x) \max(0, q(x)) \le \|p\|_\infty \max(0, q(x)).
\end{equation}
Integrating both sides over $[-2, 2]$ yields $1 \le \|p\|_\infty \int_{-2}^2 \max(0, q(x)) \, dx$. Dividing by $D = \int_{-2}^2 \max(0, q(x)) \, dx > 0$ yields the result.
\end{proof}

\section{Exact Fourth-Moment Identity}
\label{sec:fourth}

Moments of the overlap function were used by Moser and Murdeshwar \cite{mosermurdeshwar1966} and by White \cite{white2022}. For any $f \in \mathcal{H}$, let $F_m = \int_{-1}^1 t^m f(t) \, dt$ and $G_m = \int_{-1}^1 t^m g(t) \, dt$. Since $f(t) + g(t) = 1$:
\begin{equation}
F_m + G_m = \int_{-1}^1 t^m \, dt = \begin{cases} \frac{2}{m+1}, & m \text{ even}, \\ 0, & m \text{ odd}. \end{cases}
\end{equation}
In particular, $F_0 = G_0 = 1$, $F_1 + G_1 = 0$, and $F_2 + G_2 = 2/3$. The second moment of $M(x)$ satisfies:
\begin{equation}
\int_{-2}^2 x^2 M(x) \, dx = G_2 F_0 - 2G_1F_1 + G_0F_2 = \left(\frac{2}{3} - F_2\right) + 2F_1^2 + F_2 = \frac{2}{3} + \frac{1}{2}\E(M)^2,
\end{equation}
where $\E(M) = G_1 - F_1 = -2F_1$. Thus $F_2$ cancels identically from the second moment. We extend this cancellation to the fourth moment:

\begin{theorem}[Fourth-Moment Balance Identity]
\label{thm:fourth_moment}
For any $f \in \mathcal{H}$ with $M(x) = \int f(t)g(x+t) \, dt$, the fourth moment satisfies:
\begin{equation}
\int_{-2}^2 x^4 M(x) \, dx = \frac{2}{5} + 4F_2 - 6F_2^2 + 8F_1F_3,
\end{equation}
where $F_m = \int_{-1}^1 t^m f(t) \, dt$. In particular, $F_4$ cancels identically, and:
\begin{equation}
\int_{-2}^2 x^4 M(x) \, dx \le \frac{16}{15} + |\E(M)|.
\end{equation}
\end{theorem}

\begin{proof}
Expanding $(y - t)^4 = y^4 - 4y^3t + 6y^2t^2 - 4yt^3 + t^4$:
\begin{equation}
\int_{-2}^2 x^4 M(x) \, dx = G_4 F_0 - 4 G_3 F_1 + 6 G_2 F_2 - 4 G_1 F_3 + G_0 F_4.
\end{equation}
Substituting $G_4 = 2/5 - F_4$, $G_3 = -F_3$, $G_2 = 2/3 - F_2$, $G_1 = -F_1$, and $F_0 = G_0 = 1$:
\begin{equation}
\int_{-2}^2 x^4 M(x) \, dx = \left(\frac{2}{5} - F_4\right) + 4F_1F_3 + 6\left(\frac{2}{3} - F_2\right)F_2 + 4F_1F_3 + F_4.
\end{equation}
The quartic moments $+F_4$ and $-F_4$ cancel, yielding $\frac{2}{5} + 4F_2 - 6F_2^2 + 8F_1F_3$. Completing the square in $F_2$:
\begin{equation}
\frac{2}{5} + 4F_2 - 6F_2^2 = \frac{16}{15} - 6\left(F_2 - \frac{1}{3}\right)^2 \le \frac{16}{15}.
\end{equation}
By rearrangement, $|F_3| \le \int_0^1 t^3 \, dt = 1/4$.
Because $\E(M) = -2F_1$, we have $|8F_1F_3| = 4|\E(M)| \, |F_3| \le |\E(M)|$,
which yields the bound $\frac{16}{15} + |\E(M)|$.
\end{proof}

\begin{remark}
The certificate of Section~\ref{sec:certificate} uses only the moment functions $x$ and $x^2$, through the exact identity $\int x^2 M = 2/3 + \E(M)^2/2$. The bound of Theorem~\ref{thm:fourth_moment} gives an admissible constraint for $\phi(x) = x^4$, which is not used there.
\end{remark}

\section{Certified Continuous Lower Bound: \texorpdfstring{$\mu > 0.38056171$}{mu > 0.38056171}}
\label{sec:certificate}

\subsection{Mean-Conditioned Bins and Central Subdivision}

Following Price \cite{price2026}, the range of the mean $\bar{x} = \E(M) = \int_{-2}^2 x M(x) \, dx$ is partitioned into bins $[a_k, b_k]$; since $\E(M) = -2F_1$ and $|F_1| \le 1/2$, these bins cover $[-1, 1]$. For each bin, nonnegative multipliers $\lambda_j$ define the dual function $q_k(x) = 1 + \sum_j \lambda_j (B_j - \phi_j(x))$, where each constraint $\int_{-2}^2 M_f \phi_j \le B_j$ holds for every $f \in \mathcal{H}$ whose mean lies in the bin. By Proposition~\ref{thm:majorant_duality}, every such $f$ satisfies $\|M_f\|_\infty \ge 1/D_k$, where $D_k = \int_{-2}^2 \max(0, q_k(x)) \, dx$, and therefore $\mu \ge 1/\max_k D_k$.

The binding bins are those adjacent to $\bar{x} = 0$, where the variance $2/3 - \bar{x}^2/2$ of $M$ is largest. In the certificates of Price and Deng the two central bins are $[-1/320, 0]$ and $[0, 1/320]$, so that the second-moment bound used there is $2/3 + 1/204800$. The certificate of this paper (\nolinkurl{certs/erdos_elevated_candidate.json}) has 178 bins. The interval $[-1/64, 1/64]$ is divided into 16 bins of width $1/512$; in particular the central bins are $[-1/512, 0]$ and $[0, 1/512]$, for which
\begin{equation}
B(x^2) = \frac{2}{3} + \frac{1}{2}\left(\frac{1}{512}\right)^2 = \frac{2}{3} + \frac{1}{524288}.
\end{equation}
The remaining 162 bins cover $[-1, -1/64] \cup [1/64, 1]$.

\subsection{Test Function Families and Analytical Upper Bounds}
\label{subsec:families}

For each bin $[a_k, b_k]$ the following bounds hold for every $f \in \mathcal{H}$ with mean in the bin. The verifier recomputes each $B_j$ from its definition and does not use the stored values.
\begin{enumerate}
\item \textbf{Mean Constraints ($\phi(x) = \pm x$)}:
  $\int_{-2}^2 x M(x) \, dx \le b_k$ and $\int_{-2}^2 (-x) M(x) \, dx \le -a_k$.
\item \textbf{Second Moment Constraints ($\phi(x) = \pm x^2$)}:
  By the identity $\int x^2 M(x) \, dx = 2/3 + \bar{x}^2/2$ of Section~\ref{sec:fourth}, the bound for $\phi(x) = x^2$ is $B(x^2) = 2/3 + \max(a_k^2, b_k^2)/2$, and the bound for $\phi(x) = -x^2$ is $-2/3 - \min_{\bar{x} \in [a_k, b_k]} \bar{x}^2/2$.
\item \textbf{Weighted Parseval Constraint}:
  For integer harmonics $\xi = n\pi$ ($n \in \N$), $\operatorname{sinc}(n\pi) = 0$, which forces $C_{n\pi} = \int_{-2}^2 \cos(n\pi x) M(x) \, dx = -|\widehat{f}(n\pi)|^2 \le 0$. By Parseval's identity applied to the centered weight $h(t) = f(t) - 1/2$, the Fourier coefficients satisfy $\sum_{n=1}^\infty |\widehat{f}(n\pi)|^2 = \int_{-1}^1 h(t)^2 \, dt \le 1/2$. Therefore, for any weights $(w_n)_{n=1}^{400}$ satisfying $0 \le w_n \le 1$ (see \cite{white2022, deng2026}),
  \begin{equation}
  -\sum_{n=1}^{400} w_n \int_{-2}^2 \cos(n\pi x) M_f(x) \, dx = \sum_{n=1}^{400} w_n |\widehat{f}(n\pi)|^2 \le \frac{1}{2}.
  \end{equation}
  In a dual function this constraint enters with a multiplier $\lambda_W > 0$, together with the constraints $\int \cos(n\pi x) M_f(x) \, dx \le 0$ with multipliers $\eta_n \ge 0$, so that $\cos(n\pi x)$ receives the combined coefficient $c_n = \lambda_W w_n - \eta_n$. Following Deng \cite{deng2026}, the certificate stores $\lambda_W$ and the $c_n$; any $c_n \le \lambda_W$ admits such a decomposition, with $w_n = c_n/\lambda_W$ and $\eta_n = 0$ if $c_n \ge 0$, and with $w_n = 0$ and $\eta_n = -c_n$ otherwise. In the certificate this constraint appears in the 16 bins of width $1/512$, with $n \le 400$, and all 6400 stored coefficients satisfy $c_n \le \lambda_W$ in exact decimal arithmetic, without a feasibility tolerance.
\item \textbf{Trigonometric Constraints from the Complex Lens}:
  The cosine, sine and line-bundle bounds follow from an enclosure of the Fourier transform:
  \begin{lemma}[Complex Lens Bounds for Fourier Constraints]
  \label{lem:lens}
  For any $f \in \mathcal{H}$ and $\xi \in \R$, let $s = \operatorname{sinc}(\xi)$, with $\operatorname{sinc}(0)=1$, and write $\widehat{f}(\xi) = u + iv$. Then $\widehat{f}(\xi)$ lies within the closed complex lens:
  \begin{equation}
  \mathcal{L}_s = \left\{ (u, v) \in \R^2 : u^2 + v^2 \le 1 \quad \text{and} \quad (u - 2s)^2 + v^2 \le 1 \right\}.
  \end{equation}
  Consequently, the overlap function $M_f(x)$ satisfies the following analytical upper bounds for all $\xi \in \R$ and $\alpha \in \R$:
  \begin{align}
  \int_{-2}^2 \cos(\xi x) M_f(x) \, dx &\le s^2 = \operatorname{sinc}^2(\xi), \label{eq:bound_cos} \\
  -\int_{-2}^2 \cos(\xi x) M_f(x) \, dx &\le 1 - 2s^2 = 1 - 2\operatorname{sinc}^2(\xi), \label{eq:bound_negcos} \\
  \left| \int_{-2}^2 \sin(\xi x) M_f(x) \, dx \right| &\le 2 |s| \sqrt{1 - s^2} = 2 |\operatorname{sinc}(\xi)| \sqrt{1 - \operatorname{sinc}^2(\xi)}, \label{eq:bound_sin} \\
  \int_{-2}^2 (\cos(\xi x) + \alpha \sin(\xi x)) M_f(x) \, dx &\le (1 + \alpha^2) s^2 = (1 + \alpha^2) \operatorname{sinc}^2(\xi). \label{eq:bound_linebundle}
  \end{align}
  In particular, if $\xi = k\pi$ with $k \in \Z \setminus \{0\}$, then $s = 0$, and \eqref{eq:bound_sin} gives $\int_{-2}^2 \sin(k\pi x) M_f(x) \, dx = 0$.
  \end{lemma}
  \begin{proof}
  Since $0 \le f(t) \le 1$ with $\supp(f) \subseteq [-1, 1]$ and $\int_{-1}^1 f(t) \, dt = 1$, we have $|\widehat{f}(\xi)| \le \int_{-1}^1 |f(t)| \, dt = 1 \implies u^2 + v^2 \le 1$. Similarly, $g(t) = \mathbf{1}_{[-1, 1]}(t) - f(t) \ge 0$ satisfies $\int_{-1}^1 g(t) \, dt = 1$, so $|\widehat{g}(\xi)| \le 1$. Because $\widehat{g}(\xi) = 2s - \widehat{f}(\xi) = (2s - u) - iv$, we obtain $(2s - u)^2 + v^2 \le 1$, establishing that $(u, v) \in \mathcal{L}_s$.
  
  As derived in Section~\ref{sec:continuous}, $\int_{-2}^2 \cos(\xi x) M_f(x) \, dx = 2su - u^2 - v^2$ and $\int_{-2}^2 \sin(\xi x) M_f(x) \, dx = 2sv$.
  For \eqref{eq:bound_cos}: $2su - u^2 - v^2 = s^2 - (u - s)^2 - v^2 \le s^2$. Averaging the two disk inequalities gives $(u-s)^2+v^2 \le 1-s^2$. Hence $u^2+v^2-2su \le 1-2s^2$, proving \eqref{eq:bound_negcos}; equality holds on the lens at $u=s$, $v^2=1-s^2$. The same averaged inequality gives $v^2 \le 1-s^2$. For \eqref{eq:bound_sin}: $|2sv| = 2|s||v| \le 2|s|\sqrt{1 - s^2}$. For \eqref{eq:bound_linebundle}: $\operatorname{Re}\widehat{M} - \alpha \operatorname{Im}\widehat{M} = 2su - u^2 - v^2 + 2\alpha sv = (1 + \alpha^2)s^2 - (u - s)^2 - (v - \alpha s)^2 \le (1 + \alpha^2)s^2$.
  \end{proof}
\item \textbf{Fourth moment}:
  Theorem~\ref{thm:fourth_moment} bounds the constraint $\phi(x)=x^4$. Cauchy--Schwarz and $\int M=1$ give $\int x^4M \ge (\int x^2M)^2 \ge 4/9$, so $B(-x^4)=-4/9$ is also admissible. Neither constraint is used in the present certificate.
\end{enumerate}

The certificate uses \eqref{eq:bound_cos} and \eqref{eq:bound_sin} at 156 distinct finite-decimal frequencies. None is an exact nonzero integer multiple of $\pi$; the exact integer harmonics are represented separately in the Parseval constraint. It uses \eqref{eq:bound_linebundle} in 299 instances, with 146 distinct pairs $(\xi, \alpha)$, 37 distinct frequencies and $\alpha \in \{\pm\frac{1}{8}, \pm\frac{1}{4}, \pm\frac{1}{2}, \pm 1, \pm 2, \pm 3\}$, in 22 bins. We call the constraints \eqref{eq:bound_linebundle} line-bundle constraints: for fixed $\xi$ they form a family indexed by the direction $(1, \alpha)$ in the plane spanned by $\cos(\xi x)$ and $\sin(\xi x)$, and $\cos(\xi x) + \alpha \sin(\xi x) = \sqrt{1 + \alpha^2} \cos(\xi x - \arctan \alpha)$ is a phase-shifted cosine.

\subsection{Interval Arithmetic Verification}

The verifier \nolinkurl{scripts/prove_erdos.py} extends the Arb verifier of Price \cite{price2026} to the constraint families above. It checks the exact decimal sign and Parseval coefficient conditions before integration and rejects unsupported constraints. Each worker uses its own Arb precision context. For each bin the verifier recomputes every right-hand side $B_j$ as an outward Arb enclosure, increased by $10^{-14}$. Candidate sign changes of $q_k$ on $[-2, 2]$ are located by floating-point sampling, and the resulting subdivision is refined by bisection. On each cell $[c - r, c + r]$ the verifier evaluates $q_k$ and its first three derivatives at $c$ in Arb and bounds the Taylor remainder by $M_4 r^4/24$, where $M_4$ is an upper bound for $\sup |q_k^{(4)}|$ computed from the multipliers and frequencies (option \texttt{--order 4}). A cell on which the enclosure of $q_k$ is nonpositive contributes nothing; a cell on which it is nonnegative contributes the Arb enclosure of $\int q_k$ obtained from an exact antiderivative; a cell of width at most $3 \times 10^{-5}$ on which the sign is not decided contributes its width times an upper bound for $q_k$. Every contribution is an upper bound, so the sum bounds $D_k$ from above, and floating-point arithmetic only influences the choice of cells. Two stored bin endpoints differ by $1.3 \times 10^{-16}$; the verifier closes this gap by enlarging one of the two bins, which only weakens its constraints. The 160-bit run certifies all 178 bins. The per-bin bounds and cell counts are recorded in \nolinkurl{certs/verified_elevated_arb_bins.csv}. Table~\ref{tab:certified_bins} lists the largest values.

\begin{table}[h]
\centering
\small
\begin{tabular}{lccc}
\toprule
\textbf{Bins} & \textbf{Mean interval} & \textbf{Certified $D \le$} & \textbf{$\|M_f\|_\infty \ge$} \\
\midrule
88 & $[-1/512, 0]$ & $2.62769470$ & $0.38056171$ \\
89 & $[0, 1/512]$ & $2.62769470$ & $0.38056171$ \\
87 & $[-1/256, -1/512]$ & $2.62767321$ & $0.38056482$ \\
90 & $[1/512, 1/256]$ & $2.62767321$ & $0.38056482$ \\
remaining 174 bins & $|\bar{x}| \ge 1/256$ & $2.62753852$ & $0.38058433$ \\
\bottomrule
\end{tabular}
\caption{Upper bounds for $D_k = \int_{-2}^2 \max(0,q_k(x))\,dx$ from the 160-bit Arb verification. Bins are numbered from $0$ as in \nolinkurl{certs/erdos_elevated_candidate.json}. The last column is conditional on the stated mean interval and uses the reciprocal of the displayed upper bound. The final row bounds all remaining bins.}
\label{tab:certified_bins}
\end{table}

\begin{theorem}
\label{thm:main}
The minimum overlap constant satisfies
\begin{equation}
\mu \ge \frac{1}{D_{\max}} > 0.38056171, \qquad D_{\max} \le 2.62769470.
\end{equation}
\end{theorem}
\begin{proof}
By \eqref{eq:sd} it suffices to bound $\|M_f\|_\infty$ uniformly below by a constant greater than $0.38056171$. The mean of $M_f$ lies in $[-1, 1]$, which is covered by the 178 bins of the certificate. In each bin every constraint is one of the admissible inequalities of Section~\ref{subsec:families} with a nonnegative multiplier, so Proposition~\ref{thm:majorant_duality} applies with $p = M_f$. The Arb verification gives $D_k \le 2.62769470$ for every bin $k$, with the largest computed upper bound in the bin $[-1/512,0]$. Hence $\|M_f\|_\infty \ge 1/2.62769470 > 0.38056171$, uniformly in $f$.
\end{proof}

\begin{remark}[Independent check of the central window]
\label{rem:mpfr}
The C program \nolinkurl{scripts/verify_quadrisected_mpfr.c} is derived from Deng's MPFR verifier \cite{deng2026}, with the second-moment bound parameterized as $2/3+1/(2d^2)$ for the window $|\bar{x}|\le 1/d$. The default is $d=1280$. Before compilation, \nolinkurl{scripts/verify_central_mpfr.py} checks multiplier signs, the exact harmonic index set $1,\ldots,400$, the inequalities $c_n\le\lambda_W$, and that the manifest is a gap-free partition of $[0,2]$. It evaluates Deng's central certificate (\nolinkurl{scripts/weighted_center_combined_038056070.txt}: the second-moment multiplier, 80 cosine constraints and 400 combined harmonic coefficients, all even in $x$) with 192-bit directed rounding and a sixth-order Taylor enclosure with a global bound on the seventh derivative, over 32 chunks of $[0, 2]$, and uses $D = 2\int_0^2 \max(0, q)$. Since the certificate has no odd terms, it applies to every $f$ with $|\bar{x}| \le m$ once the second-moment bound for $m$ is used. Table~\ref{tab:window} lists the resulting bounds for several values of $m$; the driver selects each window with \texttt{--mean-denominator}. For the window $|\bar{x}| \le 1/512$ the certificate of this paper gives the smaller value $2.62769470$, so the additional constraint families account for part of the improvement over $0.38056070$, beyond the narrower central bins. Theorem~\ref{thm:main} rests on the Arb verification alone; the MPFR computation uses a different program and certificate on a smaller mean window, not a second verification of the 178-bin certificate.
\end{remark}

\begin{table}[h]
\centering
\small
\begin{tabular}{cccc}
\toprule
\textbf{Window} & \textbf{$B(x^2) - 2/3$} & \textbf{Certified $D \le$} & \textbf{$\|M_f\|_\infty \ge$} \\
\midrule
$|\bar{x}| \le 1/320$ & $1/204800$ & $2.62770157$ & $0.38056072$ \\
$|\bar{x}| \le 1/512$ & $1/524288$ & $2.62769646$ & $0.38056146$ \\
$|\bar{x}| \le 1/1280$ & $1/3276800$ & $2.62769373$ & $0.38056185$ \\
$|\bar{x}| \le 1/5120$ & $1/52428800$ & $2.62769324$ & $0.38056192$ \\
\bottomrule
\end{tabular}
\caption{Deng's central certificate evaluated by the 192-bit MPFR verifier with the second-moment bound of a central window $|\bar{x}| \le m$. The last column applies only within the stated window.}
\label{tab:window}
\end{table}

\subsection{Trust Model}

Linear programming (HiGHS through SciPy) is used only to produce candidate multipliers $\lambda_j \ge 0$ and frequencies. The proof of Theorem~\ref{thm:main} depends on the validity of each constraint family (Section~\ref{subsec:families}), on Proposition~\ref{thm:majorant_duality}, and on the computation of upper bounds for $D_k$, which uses Arb ball arithmetic with outward rounding together with the conservative rules for undecided cells described above.

\section{Numerical Observations on Discretized Profiles}

\subsection{Discrete Cross-Correlation Mass Conservation}

For a discrete board vector $v \in [0, 1]^N$ with $\sum_{i=1}^N v_i = N/2$, define the normalized cross-correlation at lag $m \in [-(N-1), N-1]$ by:
\begin{equation}
S(m, v) = \frac{2}{N} \sum_{i = \max(1, 1-m)}^{\min(N, N-m)} v_i (1 - v_{i+m}).
\end{equation}
\begin{proposition}[Discrete Cross-Correlation Mass Conservation]
\label{prop:discrete_mass}
For any vector $v \in [0, 1]^N$ with $\sum_{i=1}^N v_i = N/2$:
\begin{equation}
\sum_{m = -(N-1)}^{N-1} S(m, v) = \frac{N}{2}.
\end{equation}
\end{proposition}
\begin{proof}
Summing $S(m, v)$ over all $m$ expands to $\frac{2}{N} \sum_{i=1}^N \sum_{j=1}^N v_i(1 - v_j) = \frac{2}{N} (\sum v_i)(N - \sum v_j) = \frac{2}{N} (N/2)(N/2) = N/2$.
\end{proof}

\begin{remark}
Proposition \ref{prop:discrete_mass} represents the exact discrete counterpart of total continuous area conservation $\int_{-2}^2 M_f(x) \, dx = 1$. When averaged over the $2N - 1$ total lag channels, the mean overlap is $\frac{N/2}{2N - 1} \to \frac{1}{4} = 0.25$. At lag $m = 0$, $S(0, v) = \frac{2}{N} \sum_{i=1}^N v_i(1 - v_i)$. For any binary vector $v \in \{0, 1\}^N$, the coordinates satisfy $v_i(1 - v_i) = 0$, so $S(0, v) = 0$, meaning that the entire total mass of $N/2$ is distributed strictly over the $2N - 2$ non-zero lags.
\end{remark}

\subsection{An Exact Check of a Discretized Upper-Bound Profile}
\label{subsec:profile}

The file \nolinkurl{certs/ultra_dense_8192_best.json} contains 8192 JSON numbers, each interpreted as its parsed IEEE 754 binary64 value and then treated as an exact dyadic rational. The resulting vector $v \in [0,1]^{8192}$ is constant on consecutive blocks of 16 entries, so it describes a step function $f$ with 512 equal steps on $[-1, 1]$. For such $f$ the function $M_f$ is piecewise linear with breakpoints at multiples of the step width, so $\|M_f\|_\infty$ is a maximum over finitely many lags. The stored entries sum to slightly more than $4096$; after the uniform rescaling to mass exactly $4096$, which keeps every entry in $[0, 1]$, an exact rational evaluation gives
\begin{equation}
\mu \le \|M_f\|_\infty = 0.38085905681456\ldots < 0.38085906,
\end{equation}
with the maximum attained at the lag $x = 25/64$. The script \nolinkurl{scripts/certify_leaders.py} detects the constant blocks and performs the same evaluation by exact integer correlation on the 512-step representation. This coarsening preserves the function and its continuous overlap norm, since the overlap is affine between consecutive coarse-grid lags. The value agrees with the score $0.3808590568$ of the EinsteinArena entry of \texttt{lnzwz\_AI4M\_Agent} \cite{erdosproblems36}, and the profile does not improve the upper bounds listed in Section~1.

\section{Variational Dynamics and Equioscillation Heuristics}
\label{sec:variational}

The decomposition, parity lemma, first variation and variational identity below follow directly from the definitions. The multiplier condition is formal, and the Duffing model is heuristic. None is used in the proof of Theorem~\ref{thm:main}.

\subsection{Centered Decomposition and Parity}

For $f \in \mathcal{H}$ write $f = \frac{1}{2} + h$ on $[-1, 1]$, where $h: [-1, 1] \to [-\frac{1}{2}, \frac{1}{2}]$ satisfies $\int_{-1}^1 h(t) \, dt = 0$, and extend $h$ by zero outside $[-1, 1]$. Let $R_h(x) = \int_{\R} h(t) h(t+x) \, dt$ be the autocorrelation of $h$, and let
\begin{equation}
B_h(x) = \frac{1}{2} \int_{-1}^1 \mathbf{1}_{[-1, 1]}(x+t) \, (h(t) - h(t+x)) \, dt.
\end{equation}
Explicitly,
\begin{equation}
B_h(x) = \begin{cases} \frac{1}{2}\left(\int_{-1}^{-1+x} h(t) \, dt - \int_{1-x}^1 h(t) \, dt\right), & 0 \le x \le 2, \\[1ex] \frac{1}{2}\left(\int_{1+x}^1 h(t) \, dt - \int_{-1}^{-1-x} h(t) \, dt\right), & -2 \le x < 0. \end{cases}
\end{equation}
Expanding $f(t) g(t+x) = (\frac{1}{2} + h(t))(\frac{1}{2} - h(t+x))$ over the set of $t$ with $t, t + x \in [-1, 1]$, which has length $2 - |x|$, gives
\begin{equation}
\label{eq:centered}
M_f(x) = \frac{2 - |x|}{4} + B_h(x) - R_h(x), \qquad |x| \le 2.
\end{equation}

\begin{lemma}[Parity cancellation of the boundary term]
\label{lem:parity}
If $h(-t) = h(t)$ for all $t \in [-1, 1]$, then $B_h(x) = 0$ for all $x \in [-2, 2]$, and hence $M_f(x) = \frac{2 - |x|}{4} - R_h(x)$.
\end{lemma}
\begin{proof}
For $x \ge 0$, the change of variables $t = -v$ gives $\int_{-1}^{-1+x} h(t) \, dt = \int_{1-x}^1 h(-v) \, dv = \int_{1-x}^1 h(v) \, dv$, so the two integrals in $B_h(x)$ cancel. The case $x < 0$ is identical.
\end{proof}

Evenness of extremal profiles has not been established, and the profile of Section~\ref{subsec:profile} is not even. The rest of this section concerns even profiles.

\subsection{Formal First-Order Conditions}

For $\varphi$ supported in $[-1, 1]$ with $\int \varphi = 0$, the first variation of $M_f(x)$ in the direction $\varphi$ is $\int \varphi(t) \, (g(t+x) - f(t-x)) \, dt$, with $f$ and $g$ extended by zero. Let $f^* = \frac{1}{2} + h^*$ be an even minimizer, write $\mu^* = \|M_{f^*}\|_\infty$, and let $E = \{x \in [-2, 2] : M_{f^*}(x) = \mu^*\}$. A formal Lagrange multiplier argument suggests that there are an even probability measure $\nu$ on $E$ and a constant $\lambda$ such that
\begin{equation}
\label{eq:el_convolution}
(K_\nu h^*)(t) := \int_E h^*(t - x) \, d\nu(x) = \lambda \quad \text{for almost every } t \in \Omega = \{t \in [-1, 1] : |h^*(t)| < \tfrac{1}{2}\},
\end{equation}
with one-sided inequalities on the sets where $h^* = \pm \frac{1}{2}$. Here the constant terms cancel because $\nu$ is even. We have not justified \eqref{eq:el_convolution} rigorously; the existence of minimizers and the precise form of the multiplier rule require further work.

\begin{remark}[Variational identity]
For every even probability measure $\nu$ supported on $E$, integrating the identity $M_{f^*}(x) = \mu^*$ against $\nu$ and using Lemma~\ref{lem:parity} and Fubini's theorem gives
\begin{equation}
\mu^* = \frac{1}{2} - \frac{1}{4} \E_{\nu}[|x|] - \langle h^*, K_{\nu} h^* \rangle_{L^2([-1, 1])}.
\end{equation}
If \eqref{eq:el_convolution} held on a set of full measure, then $\langle h^*, K_\nu h^* \rangle = \lambda \int h^* = 0$ and $\mu^* = \frac{1}{2} - \frac{1}{4}\E_\nu[|x|]$. Numerically computed near-extremal step functions take the values $0$ and $1$ on sets of positive measure, so this term need not vanish.
\end{remark}

\subsection{A Duffing-Type Caricature}

Condition \eqref{eq:el_convolution} is non-local, and the constraint $|h| \le \frac{1}{2}$ is rigid. As a heuristic caricature only, replace the convolution operator by a local second-order operator and the constraint by a smooth confining nonlinearity; this leads to equations of Duffing type,
\begin{equation}
h''(t) + \alpha h(t) + \beta h(t)^3 = 0.
\end{equation}
For $\omega > 0$ and $0 < k < 1$, the function $h(t) = \frac{1}{2} \operatorname{cn}(\omega t, k)$ is an even periodic solution with $h(0) = \frac{1}{2}$ when $\alpha = \omega^2(1 - 2k^2)$ and $\beta = 8 k^2 \omega^2$; this follows from the identity $(y')^2 = \omega^2 (1 - y^2)(1 - k^2 + k^2 y^2)$ satisfied by $y(t) = \operatorname{cn}(\omega t, k)$. Its extrema lie at $\omega t \in 2K(k)\Z$, so prescribing 11 extrema on $[-1, 1]$, at $t = j/5$ for $j = -5, \dots, 5$, gives $\omega = 10 K(k)$. The modulus $k$ is not determined by this model, and the model gives no quantitative prediction for $\mu$. It is recorded only as a qualitative description of the oscillation seen in numerically computed profiles, which in addition contain plateaus where $f$ equals $0$ or $1$.

\section{Lean 4 Formalization}
\label{sec:lean}

The file \nolinkurl{formal/Erdos36.lean} is checked with Lean 4.35.0 and the Mathlib revision pinned in \nolinkurl{formal/lakefile.toml}; it contains no \texttt{sorry}. It defines $M(N)$ as in the formal-conjectures statement of Problem~36, the class of admissible weights, the overlap function $M_f$ and its moments. Its content is of two kinds.

\emph{Statements proved in Lean.} These include the polynomial bound $\frac{2}{5} + 4u - 6u^2 \le \frac{16}{15}$; the cross-term bound $8F_1F_3 \le |E|$ under $E = -2F_1$ and $|F_3| \le \frac{1}{4}$; the identity $(1+\alpha^2)s^2 - (2su - u^2 - v^2 + 2\alpha s v) = (u-s)^2 + (v-\alpha s)^2$ behind \eqref{eq:bound_linebundle}; the arithmetic step of Proposition~\ref{thm:majorant_duality}, namely that $1 \le \int pq \le \|p\|_\infty D$ and $D > 0$ imply $\|p\|_\infty \ge 1/D$; the identity $\int R_h = (\int h)^2$ conditional on explicit Fubini and translation identities, and the value $R_h(0) = \frac{1}{2}$ for $h$ with values $\pm\frac{1}{2}$; the four-term ramp formula for the length of the intersection of two intervals; and several elementary algebraic lemmas related to Section~\ref{sec:variational}.

\emph{Axioms.} The analytic inputs and the computation enter as the nine axioms listed in Table~\ref{tab:axioms}. In particular, the fourth-moment identity of Theorem~\ref{thm:fourth_moment} is stated in Lean as an axiom (its proof is the elementary computation of Section~\ref{sec:fourth}), and the statement \nolinkurl{erdos_36_lower_bound_certified}, which reads $0.38056171 \le \liminf_{N \to \infty} M(N)/N$, is derived from the axioms \nolinkurl{erdos_enhanced_arb_certificate}, \nolinkurl{erdos_certificate_bound} and \nolinkurl{continuous_reduction_to_discrete}. The integral axioms include the required integrability hypotheses, and the supremum bound assumes that $p$ is bounded above. These assumptions avoid treating nonintegrable functions as though the Bochner integral were linear on them. The Lean file therefore does not verify the numerical certificate; it records how the certificate, the duality principle and Swinnerton-Dyer's theorem combine.

\begin{table}[h]
\centering
\small
\begin{tabular}{ll}
\toprule
\textbf{Axiom} & \textbf{Content} \\
\midrule
\texttt{f\_third\_moment\_neyman\_pearson} & $|F_3| \le 1/4$ by rearrangement \\
\texttt{mean\_overlap\_fubini} & $\E(M) = G_1 - F_1$ and $F_1 + G_1 = 0$ \\
\texttt{binomial\_fourth\_moment\_identity} & the identity of Theorem~\ref{thm:fourth_moment} \\
\texttt{dual\_lower\_bound\_expansion} & linearity of $\int pq$ for integrable data \\
\texttt{pointwise\_majorant\_integral} & $\int pq \le \|p\|_\infty \int \max(0, q)$ for bounded $p \ge 0$ \\
\texttt{continuous\_reduction\_to\_discrete} & Swinnerton-Dyer's theorem, lower-bound direction \\
\texttt{erdos\_enhanced\_arb\_certificate} & the certificate exists with target $0.38056171$ \\
\texttt{erdos\_certificate\_bound} & soundness of the certificate specification \\
\texttt{erdos\_36\_upper\_bound\_construction} & $\limsup M(N)/N \le 0.38085906$ \\
\bottomrule
\end{tabular}
\caption{Axioms of \nolinkurl{formal/Erdos36.lean}.}
\label{tab:axioms}
\end{table}

\section{Computational Methodology and Reproducibility}
\label{sec:repro}

Candidate multipliers were computed with linear programming (SciPy with the HiGHS and interior-point solvers); these computations are not part of the trusted path. The repository \url{https://github.com/Destructive-Entity/erdos-minimum-overlap} contains the following files.
\begin{description}\raggedright
\item[{\nolinkurl{certs/erdos_elevated_candidate.json}}] The 178-bin dual certificate.
\item[{\nolinkurl{scripts/prove_erdos.py}}] The Arb verifier, which extends the verifier of Price \cite{price2026}; \nolinkurl{certs/verified_elevated_arb_report.json} is the report of a full run.
\item[Central-window check.] \nolinkurl{scripts/verify_central_mpfr.py} compiles \nolinkurl{scripts/verify_quadrisected_mpfr.c}, which is Deng's MPFR verifier with the modified second-moment constant, and runs it over the 32 chunks listed in \nolinkurl{scripts/weighted_038056070_center_manifest.csv}. The coefficient file \nolinkurl{scripts/weighted_center_combined_038056070.txt} is Deng's central certificate \cite{deng2026}, unchanged.
\item[Discrete profile.] \nolinkurl{certs/ultra_dense_8192_best.json} contains the profile of Section~\ref{subsec:profile}; \nolinkurl{scripts/certify_leaders.py} evaluates it.
\item[{\nolinkurl{formal/Erdos36.lean}}] The Lean file of Section~\ref{sec:lean}.
\end{description}
The aggregate Arb report records the certificate and verifier SHA-256 hashes; \nolinkurl{certs/verified_elevated_arb_bins.csv} records each bin. Regression tests compare derivatives with the antiderivative and check rejection of inadmissible inputs, incomplete mean covers and incomplete MPFR partitions. The commands are listed in Appendix~\ref{app:code}.

\paragraph{AI assistance.}
Kreflux AI's internally developed models assisted manuscript preparation and computational-code development and review. The author retains full responsibility for all mathematical claims, code, numerical results and final text.

\begin{thebibliography}{99}\sloppy
\bibitem{erdos1955} P.~Erd\H{o}s, \emph{Some remarks on number theory}, Riveon Lematematika \textbf{9} (1955), 45--48 (in Hebrew).
\bibitem{motzkin1956} T.~S.~Motzkin, K.~E.~Ralston, and J.~L.~Selfridge, \emph{Minimal overlapping under translation}, Bull. Amer. Math. Soc. \textbf{62} (1956), 558 (abstract).
\bibitem{moser1959} L.~Moser, \emph{On the minimal overlap problem of Erd\H{o}s}, Acta Arith. \textbf{5} (1959), 117--119, doi:10.4064/aa-5-2-117-119.
\bibitem{mosermurdeshwar1966} L.~Moser and M.~G.~Murdeshwar, \emph{On the overlap of a function with the translation of its complement}, Colloq. Math. \textbf{15} (1966), 93--97.
\bibitem{haugland1996} J.~K.~Haugland, \emph{Advances in the minimum overlap problem}, J. Number Theory \textbf{58} (1996), 71--78.
\bibitem{haugland2016} J.~K.~Haugland, \emph{The minimum overlap problem revisited}, arXiv:1609.08000 (2016).
\bibitem{guy2004} R.~K.~Guy, \emph{Unsolved Problems in Number Theory}, 3rd ed., Springer, New York, 2004, Problem C17.
\bibitem{erdosproblems36} T.~F.~Bloom, \emph{Erd\H{o}s Problems}, Problem \#36 and its discussion thread, \url{https://www.erdosproblems.com/36} (accessed November 2026).
\bibitem{white2022} E.~P.~White, \emph{Erd\H{o}s' minimum overlap problem}, arXiv:2201.05704 (2022).
\bibitem{white2023} E.~P.~White, \emph{A new bound for Erd\H{o}s' minimum overlap problem}, Acta Arith. \textbf{208} (2023), 235--255.
\bibitem{kimpilanci2026} S.~Kim and M.~Pilanci, \emph{AI-assisted discovery of convex relaxations via dual agents}, arXiv:2606.31182 (2026).
\bibitem{price2026} L.~Price, \emph{A 0.38055470 lower bound for Erd\H{o}s Problem 36}, proof note and certificate package, \url{https://github.com/Leeham06972452/erdos-36-lower-bound} (2026).
\bibitem{raval2026} R.~R.~Raval, \emph{The Erd\H{o}s minimum-overlap constant exceeds 0.38055925}, version 0.3.1, Zenodo, \url{https://doi.org/10.5281/zenodo.22647743} (2026); code at \url{https://github.com/ruturajr-raval/erdos-minimum-overlap-lower-bound}.
\bibitem{deng2026} H.~Deng, \emph{A sharper certified lower bound for Erd\H{o}s' minimum-overlap problem}, version 1.1.1, Zenodo, \url{https://doi.org/10.5281/zenodo.22279894} (2026); code at \url{https://github.com/emerardd/erdos-36-lower-bound}.
\bibitem{machillari2026} B.~Machillari Dowgu\~{n}iec, \emph{A certified lower bound for Erd\H{o}s's 36 (minimum-overlap problem) via subadditivity constraints and exact arithmetic}, manuscript, 19 September 2026, \url{https://drive.google.com/file/d/1ThaY0esdiK47E6jL2IGkh11UxkeqzRPy/view}.
\bibitem{budala2026} C.~Budala, \emph{Explicit bounds for Erd\H{o}s's minimum-overlap problem}, preprint with exact-arithmetic checkers, \url{https://github.com/CristianBudala/erdos-36-bounds} (2026).
\bibitem{alphaevolve2025} B.~Georgiev, J.~G\'omez-Serrano, T.~Tao, and A.~Z.~Wagner, \emph{Mathematical exploration and discovery at scale}, arXiv:2511.02864 (2025).
\bibitem{tttdiscover2026} M.~Yuksekgonul, D.~Koceja, X.~Li, F.~Bianchi, J.~McCaleb, X.~Wang, J.~Kautz, Y.~Choi, J.~Zou, C.~Guestrin, and Y.~Sun, \emph{Learning to discover at test time}, arXiv:2601.16175 (2026).
\bibitem{simpletes2026} H.~Ye, H.~Lin, J.~Tang, Y.~Luo, C.~Yang, C.~Su, R.~Thapa, R.~Yang, R.~Liu, Z.~Li, et al., \emph{Evaluation-driven scaling for scientific discovery}, arXiv:2604.19341 (2026).
\bibitem{einsteinarena2026} F.~Bianchi, Y.~Kwon, A.~Pappu, and J.~Zou, \emph{Harnessing the collective intelligence of AI agents in the wild for new discoveries}, arXiv:2606.10402 (2026); leaderboard at \url{https://einsteinarena.com}.
\end{thebibliography}

\appendix

\section{Verification Commands}
\label{app:code}

From the repository root, with the Python packages \texttt{python-flint}, \texttt{mpmath}, \texttt{numpy} and \texttt{scipy}, a C compiler, and the GNU MPFR and GMP libraries:
\begin{verbatim}
python3 scripts/prove_erdos.py certs/erdos_elevated_candidate.json \
    --order 4 --prec 160 --target 0.38056171 --jobs 2
python3 scripts/verify_central_mpfr.py
python3 scripts/certify_leaders.py certs/ultra_dense_8192_best.json
\end{verbatim}
The first command certifies every bin against the target $0.38056171$ and reports the worst bin. The second compiles the MPFR verifier and checks the window $|\bar{x}| \le 1/1280$, reporting $D \le 2.62769373$. The third evaluates the profile of Section~\ref{subsec:profile} by exact integer correlation after lossless block compression. To reproduce every row of Table~\ref{tab:window}, pass \texttt{--mean-denominator} with values $320$, $512$, $1280$ and $5120$ to the second command. Run the regression tests with \texttt{python3 -m unittest discover -s scripts}.

\end{document}
